\section{Threat Hunting Playbooks}

Threat hunting playbooks provide structured, repeatable approaches for investigating specific adversary behaviors. Unlike detections, which run continuously, playbooks guide analysts through deeper manual exploration. Each playbook begins with a hypothesis, identifies relevant telemetry, and provides SPL, KQL, and Linux commands that support the investigation. These playbooks are designed to be practical, adaptable, and aligned with real-world adversary techniques.

\subsection*{Playbook: Suspicious PowerShell Activity}

\textbf{Hypothesis:} An adversary is using PowerShell to execute encoded or obfuscated commands.

\textbf{Relevant Telemetry:} Sysmon Event ID 1, PowerShell logs (4103, 4104), process command-line arguments.

\textbf{SPL Example}
\begin{lstlisting}
index=sysmon EventCode=1 Image="*powershell.exe"
| search CommandLine="*EncodedCommand*"
| table _time, Computer, User, ParentImage, Image, CommandLine
\end{lstlisting}

\textbf{KQL Equivalent}
\begin{lstlisting}
Sysmon
| where EventID == 1 and Image endswith "powershell.exe"
| where CommandLine contains "EncodedCommand"
| project TimeGenerated, Computer, User, ParentImage, Image, CommandLine
\end{lstlisting}

\textbf{Linux Pivot (Sysmon for Linux)}
\begin{lstlisting}[language=bash]
grep -i "powershell" /var/log/sysmon/sysmon.log
\end{lstlisting}

This playbook helps identify obfuscated execution, a common technique in malware delivery and post-exploitation.

\subsection*{Playbook: Credential Access Attempts}

\textbf{Hypothesis:} An adversary is attempting to brute-force or misuse credentials.

\textbf{Relevant Telemetry:} Windows Event ID 4625, 4624, Linux auth logs, network logs.

\textbf{SPL Example}
\begin{lstlisting}
index=wineventlog EventCode=4625
| stats count by AccountName, WorkstationName
| where count > 10
\end{lstlisting}

\textbf{KQL Equivalent}
\begin{lstlisting}
SecurityEvent
| where EventID == 4625
| summarize Count=count() by Account, Workstation
| where Count > 10
\end{lstlisting}

\textbf{Linux Pivot}
\begin{lstlisting}[language=bash]
grep "Failed password" /var/log/auth.log | awk '{print $9}' | sort | uniq -c
\end{lstlisting}

This playbook identifies brute-force attempts, password spraying, and compromised accounts.

\subsection*{Playbook: Lateral Movement via Remote Execution}

\textbf{Hypothesis:} An adversary is moving laterally using remote execution tools such as PsExec, WMI, or SSH.

\textbf{Relevant Telemetry:} Sysmon Event ID 1, 3, Windows Event ID 7045, Linux SSH logs.

\textbf{SPL Example (PsExec)}
\begin{lstlisting}
index=sysmon EventCode=1 Image="*psexec*"
| table _time, Computer, User, ParentImage, Image, CommandLine
\end{lstlisting}

\textbf{KQL Equivalent}
\begin{lstlisting}
Sysmon
| where EventID == 1 and Image contains "psexec"
| project TimeGenerated, Computer, User, ParentImage, Image, CommandLine
\end{lstlisting}

\textbf{Linux Pivot (SSH)}
\begin{lstlisting}[language=bash]
grep "Accepted password" /var/log/auth.log
grep "Failed password" /var/log/auth.log
\end{lstlisting}

This playbook helps identify unauthorized remote execution and lateral movement.

\subsection*{Playbook: Rare Network Connections}

\textbf{Hypothesis:} A compromised host is communicating with an unusual or suspicious external destination.

\textbf{Relevant Telemetry:} Sysmon Event ID 3, firewall logs, DNS logs.

\textbf{SPL Example}
\begin{lstlisting}
index=sysmon EventCode=3
| stats count by DestinationIp, DestinationPort
| where count < 5
\end{lstlisting}

\textbf{KQL Equivalent}
\begin{lstlisting}
Sysmon
| where EventID == 3
| summarize Count=count() by DestinationIp, DestinationPort
| where Count < 5
\end{lstlisting}

\textbf{Linux Pivot}
\begin{lstlisting}[language=bash]
ss -tulnp
netstat -antp
\end{lstlisting}

This playbook is effective for identifying command-and-control channels and exfiltration paths.

\subsection*{Playbook: Persistence Mechanisms}

\textbf{Hypothesis:} An adversary has established persistence on a host.

\textbf{Relevant Telemetry:} Scheduled tasks, services, registry modifications, Linux cron jobs.

\textbf{SPL Example (Scheduled Tasks)}
\begin{lstlisting}
index=wineventlog EventCode=4698
| table _time, TaskName, Author, Command
\end{lstlisting}

\textbf{KQL Equivalent}
\begin{lstlisting}
SecurityEvent
| where EventID == 4698
| project TimeGenerated, TaskName, Author, Command
\end{lstlisting}

\textbf{Linux Pivot (Cron Jobs)}
\begin{lstlisting}[language=bash]
cat /etc/crontab
ls -al /etc/cron.d/
\end{lstlisting}

This playbook helps identify long-term footholds established by attackers.

\subsection*{Playbook: Suspicious File Modifications}

\textbf{Hypothesis:} An adversary is modifying system files or dropping malicious payloads.

\textbf{Relevant Telemetry:} Sysmon Event ID 11, Linux file modification logs.

\textbf{SPL Example}
\begin{lstlisting}
index=sysmon EventCode=11
| table _time, Computer, User, TargetFilename
\end{lstlisting}

\textbf{KQL Equivalent}
\begin{lstlisting}
Sysmon
| where EventID == 11
| project TimeGenerated, Computer, User, TargetFilename
\end{lstlisting}

\textbf{Linux Pivot}
\begin{lstlisting}[language=bash]
find / -type f -mtime -1 2>/dev/null
\end{lstlisting}

This playbook is useful for detecting tampering, malware staging, and unauthorized changes.

\subsection*{Playbook Structure Summary}

Each playbook follows the same structure:
\begin{itemize}
    \item Hypothesis
    \item Relevant telemetry
    \item SPL query
    \item KQL equivalent
    \item Linux investigation commands
\end{itemize}

This consistent format ensures that hunts are repeatable, scalable, and aligned with real-world adversary behavior.
